\section{The optimal power}\label{sec:exponent}

We combine the quantitative bounds before recalling the universal upper
bound in the appendix.

\begin{proof}[Proof of Theorem~\ref{thm:main}]
For each fixed $k\ge2$, Proposition~\ref{prop:quantum} gives
$\Q(F_{k,m})\le C_k\sqrt m(\log m)^{b_k}$ for suitable constants.
Proposition~\ref{prop:classical} gives
\[
 \R(F_{k,m})>
 \left\lfloor\frac m{100}
  \left\lfloor\frac{m^{1-1/k}}{(\log m)^2}\right\rfloor
 \right\rfloor
 \ge c_k\frac{m^{2-1/k}}{(\log m)^2}
\]
for sufficiently large $m$, after decreasing $c_k>0$.
This proves \eqref{eq:main-family}.

Now fix $0\le\alpha<4$ and choose an integer $k\ge2$ such that
$2-1/k>\alpha/2$. The integer $k$ remains fixed as $m$ tends to infinity.
The preceding bounds imply
\[
 \frac{\R(F_{k,m})}{(1+\Q(F_{k,m}))^\alpha}
 \ge c_{k,\alpha}
 \frac{m^{2-1/k-\alpha/2}}
 {(\log m)^{2+\alpha b_k}}
 \longrightarrow\infty.
\]
Hence no exponent smaller than four belongs to the set in
\eqref{eq:alpha}. Conversely, Theorem~\ref{thm:quartic-upper} proves
$\R(f)\le\D(f)=O((1+\Q(f))^4)$ for every total Boolean function.
Therefore four does belong to that set, and its infimum is four.
\end{proof}

Separate address arrays force a classical algorithm to invest in many
columns, while quantum comparisons locate a candidate before its
certificate is checked. Quantum search then supplies the quadratic
saving in checking the local tests. The direct-product estimate prevents
a classical algorithm from bypassing the addressing cost by guessing
certificate locations.
